\documentclass[margin,line,pifont,palatino,courier]{res}

\usepackage{pifont}
\usepackage[latin1] { inputenc}
\usepackage{hyperref}

%\topmargin .5in
%\oddsidemargin -.5in
%\evensidemargin -.5in
%\textwidth=6.0in
 \textheight=9.0in
%\itemsep=0in
%\parsep=0in
\usepackage{fancyhdr}
%\topmargin=0in
%\textheight=8.5in
\pagestyle{fancy}
\renewcommand{\headrulewidth}{0pt}
\fancyhf{}
%\cfoot{\thepage}
%\lfoot{\textit{\footnotesize Research Statement}}
\rfoot{{\footnotesize Curriculum Vitae, Joel Villatoro, \thepage}}


\newenvironment{list1}{
  \begin{list}{\ding{113}}{%
      \setlength{\itemsep}{0in}
      \setlength{\parsep}{0in} \setlength{\parskip}{0in}
      \setlength{\topsep}{0in} \setlength{\partopsep}{0in}
      \setlength{\leftmargin}{0.17in}}}{\end{list}}
\newenvironment{list2}{
  \begin{list}{$\bullet$}{%
      \setlength{\itemsep}{0in}
      \setlength{\parsep}{0in} \setlength{\parskip}{0in}
      \setlength{\topsep}{0in} \setlength{\partopsep}{0in}
      \setlength{\leftmargin}{0.2in}}}{\end{list}}

\begin{document}

\name{Joel Villatoro \vspace*{.1in}}

\begin{resume}

\section{\sc Contact Information}

\vspace{.05in}
\begin{tabular}{@{}p{2.75in}p{2in}}
Department of Mathematics                        & \verb+jvillato (at) iu.edu+\\
Indiana University Bloomington                &  \\
Bloomington, IN 47405                   & \\
\end{tabular}


\section{\sc Education}

{\bf Department of Mathematics, University of Illinois}\\
\vspace*{-.1in}
\begin{list1}
\item[] Ph.D. in Mathematics, August 2018

\begin{list2}
\vspace*{.05in}
\item Dissertation Title: Stacks in Poisson Geometry
\item Advisor: Rui Fernandes
\end{list2}
\end{list1}

{\bf University of Oklahoma}\\
\vspace*{-.1in}
\begin{list1}
\item[] B.S.~in Mathematics, May 2012

\end{list1}

\section{\sc Academic Positions}
{\bf Lecturer, Indiana University Bloomington}\\
\vspace*{-.1in}
\begin{list1}
\item[] Fall 2024-Present
\end{list1}

{\bf NSF Postdoc, Washington University in St. Louis}\\
\vspace*{-.1in}
\begin{list1}
\item[] Spring 2022-Spring 2024 (Funded by MPS-Ascend grant, Award No. 2137999)
\end{list1}


{\bf Guest, Max Planck Institute for Mathematics in Bonn, Germany}\\
\vspace*{-.1in}
\begin{list1}
\item[] Fall 2021
\end{list1}

{\bf Postdoctoral Fellow, KU Leuven in Belgium}\\
\vspace*{-.1in}
\begin{list1} 
\item[] September 2018-August 2021
\end{list1}



\section{\sc Honors and Awards}

\begin{tabular}{@{}p{0.8in}p{4in}}
2021 & Awarded MPS Ascend Fellowship from the NSF \\
2017--2018 & NSF AGEP (Alliances for Graduate Education and the Professoriate) Research support\\
2016  & On UIUC's `List of Teachers Ranked as Excellent' \\
2014 & NSF AGEP Research support\\
2014       & UIUC Research Board Award\\
2013 & On UIUC's `List of Teachers Ranked as Excellent' \\
2012--2013 & UIUC GAAN (Graduate Assistance in Areas of National Need) Research Fellowship
\end{tabular}


\section{\sc Publications \&\\ Preprints}

Pooja Joshi, Joel Villatoro, On the holonomy of Lie algebroids, \textit{(preprint)} \href{https://arxiv.org/abs/2608.19399}{link}

Xiang Tang, Joel Villatoro, Simplicial sheaves of modules and Morita invariance of groupoid cohomology, \textit{(preprint)} \href{https://arxiv.org/abs/2509.07285}{link}.

Angel Roman, Joel Villatoro, Convolution algebras of double groupoids and 2-groups, SIGMA 20 (2024), 093, 26 pages \href{https://arxiv.org/pdf/2312.00341}{link}.

Joel Villatoro, On the integrability of Lie algebroids by diffeological spaces, Indagationes Mathematicae, 2025, ISSN 0019-3577, \href{https://doi.org/10.1016/j.indag.2025.07.011}{link}.

Joel Villatoro, On Sheaves of Lie-Rinehart structures, arXiv:2010.15463 \textit{(preprint)} \href{https://arxiv.org/pdf/2010.15463.pdf}{link}.

Alfonso Garmendia, Joel Villatoro, 
Integration of Singular Foliations via Paths, 
International Mathematics Research Notices, 
Volume 2022, Issue 23, December 2022, 
Pages 18401-18445. \href{https://academic.oup.com/imrn/article-abstract/2022/23/18401/6362694}{link}

J. Villatoro. Picard groups of b-symplectic manifolds. \textit{Journal of Symplectic Geometry} (2021), vol. 19, no. 3, 723-775. https://dx.doi.org/10.4310/JSG.2021.v19.n3.a6. \href{https://www.intlpress.com/site/pub/files/_fulltext/journals/jsg/2021/0019/0003/JSG-2021-0019-0003-a006.pdf}{link}

J. Villatoro, Stacks in Poisson Geometry.
(2018 PhD Thesis). arXiv:1806.01939. \href{https://arxiv.org/abs/1806.01939}{link}

J. Villatoro. Poisson manifolds and their associated stacks. \textit{Letters in Mathematical Physics}
(2018), no. 3, 897-926. https://doi.org/10.1007/s11005-017-1012-5. \href{https://link.springer.com/article/10.1007/s11005-017-1012-5}{link}


\section{\sc Teaching Experience}
\begin{list1}
\item[] \textbf{Indiana University Bloomington}
\end{list1}
\begin{tabular}{@{}p{1.2in}p{0.5in}p{3in}}
Spring & 2026 & \textbf{Brief Calculus (Instructor)} 
Introductory calculus course for non-STEM majors. \\
Spring & 2026 & \textbf{Math of Decision and Beauty (Instructor)} 
Gen Ed topics course exploring mathematical ideas in decision-making and aesthetics. \\
Fall & 2025 & \textbf{Brief Calculus (Instructor)} 
Introductory calculus course for non-STEM majors. \\
Spring & 2025 & \textbf{Calculus II (Instructor)} 
Standard undergraduate calculus. \\
Spring & 2025 & \textbf{Brief Calculus (Instructor)} 
Introductory calculus course for non-STEM majors. \\
Fall & 2024 & \textbf{Calculus I (Instructor)} 
Standard undergraduate calculus for STEM majors. \\
Fall & 2024 & \textbf{Calculus II (Instructor)} 
Standard undergraduate calculus for STEM majors.\\
\end{tabular}\\
\begin{list1}
\item[] \textbf{KU Leuven}
\end{list1}\begin{tabular}{@{}p{1.2in}p{0.5in}p{3in}}
Spring & 2021  & \textbf{Student Seminar Grader}: Evaluate and provide feedback on individual bachelor student presentations \\
Fall & 2020 & \textbf{Differential Equations(TA)}: Monitor and advise students as they work individually on assigned exercises.\\
Spring & 2020 & \textbf{Symplectic Geometry(Instructor)}: Graduate level course. Design a curriculum based on a syllabus of required topics. Give lectures, write and grade exams and homework. Switched to synchronous lectures after COVID. \\
Fall & 2019 & \textbf{Wiskunde 1(TA)}: Single variable calculus for non-mathematics majors. Monitor and advise students as they work individually on assigned exercises.\\
\end{tabular}\\
% \begin{list1}
% \item[] \textbf{University of Illinois at Urbana-Champaign}
% \end{list1}
% \begin{tabular}{@{}p{1.2in}p{0.5in}p{3in}}
% Spring & 2017 & \textbf{Calculus I(TA)}: Single variable calculus. Prepare exercises and lead students in active learning sessions. Put students into groups and provide assistance and encourage cooperation and discussion.\\
% Fall   & 2016 & \textbf{Calculus I (TA)}: See above.\\
% Fall   & 2013 & \textbf{Calculus III (TA)}: Multivariable calculus. Same class format as Calculus I.\\
% Fall   & 2012 & \textbf{Linear Algebra (Tutor)}: Work in official university tutoring room for Linear Algebra students. Give hints and advise to students who come with questions and homework. 
% \end{tabular}
\begin{list1}
\item[] \textbf{University of Illinois at Urbana-Champaign}
\end{list1}
\begin{tabular}{@{}p{1.2in}p{3.5in}}
2012--2017 & \textbf{Teaching Assistant / Tutor}: Calculus I, Calculus III, Linear Algebra.
\end{tabular}
\newpage

\section{\sc Supervision}
\begin{tabular}{@{}p{1.2in}p{3in}}
Spring 2021 &  \textbf{Bachelor project supervisor}: Mentor a group of 6 Bachelor students in studying a new topic and preparing a presentation.\\
Spring 2018 & \textbf{Illinois Geometry Lab}: graduate student assistant. Graduate student supervisor for a team of students working on computation based research project about Poisson structures on Euclidean space.\\
Fall 2017 & \textbf{Illinois Geometry Lab}: graduate student assistant. Graduate student supervisor for a team of students working on computation based research project about modeling a partial differential equation.\\

\end{tabular}


\section{\sc Diversity Equity and Inclusion}

\begin{tabular}{@{}p{1.2in}p{3in}}
Fall 2023 & \textbf{Puertas College Prep}: Organizer and mentor for Junior and Senior level highschool students from disadvantaged backgrounds. \\
Spring 2023 & \textbf{WashU Mathematics DEI Committee}: Served on the math department's committee. \\
Fall 2022 & \textbf{WashU Mathematics DEI Committee}: Served on the math department's committee. \\
Fall 2022 & \textbf{Puertas College Prep Mentor and Organizer}: Organizer and mentor for Junior and Senior level highschool students from disadvantaged backgrounds. \\
2015--2017 & \textbf{SLOAN Peer Mentor}: Work with an assigned junior graduate student from underrepresented background. Provide mentorship, advise to mentee as they adapt to life as a new graduate student.\\
Summer 2012 & \textbf{Project Transformation Camp Counselor}: Worked at an educational summer camp for low income children in the OKC metro area.
\end{tabular}

\section{\sc Organizational Experience}

\begin{tabular}{@{}p{1.2in}p{0.5in}p{3.5in}}
Summer & 2025 & \textbf{Poisson Geometry, Lie Theory and Symmetry}: 
Principal organizer for an international research conference held at Instituto Superior T\'ecnico, Lisbon. Contributed to scientific planning, speaker selection, and coordination of the event. \url{https://sites.google.com/view/ipgls2025/}(\href{https://sites.google.com/view/ipgls2025/}{link}) \\
Fall & 2023 & \textbf{UIUC--WUSTL Joint Seminar Organizer}: 
Co-organized a joint research seminar between the University of Illinois and Washington University in St. Louis. \\
Spring & 2022 & \textbf{WashU Geometry Seminar Organizer}: 
Organized weekly research seminars in geometry. \\
Fall & 2022 & \textbf{WashU Geometry Seminar Organizer}: 
Organized weekly research seminars in geometry. \\
Spring & 2021 & \textbf{AGEP-GRS Conference Steering Committee}: 
Served on the steering committee for a graduate research symposium, contributing to planning and organization. \\
Fall & 2020 & \textbf{Junior Seminar Organizer}: 
Organized and led a seminar for junior graduate students. \\
Fall--Spring & 2019--2020 & \textbf{KU Leuven Weekly Poisson Meetings Organizer}: 
Organized regular research meetings in Poisson geometry. \\
Fall--Spring & 2018--2019 & \textbf{KU Leuven Weekly Poisson Meetings Organizer}: 
Organized regular research meetings in Poisson geometry. \\
\end{tabular}



\section{\sc Service}
\begin{tabular}{@{}p{1.2in}p{3in}}
Spring 2023 & \textbf{WashU Mathematics DEI Committee}: Served on the math department's committee. \\
Fall 2022 & \textbf{WashU Mathematics DEI Committee}: Served on the math department's committee. \\
Spring 2022 & \textbf{Washington University Math Club}: Organized meetings for the undergraduate mathematics club. \\
Spring 2022 & \textbf{Glenridge Elementary School Math Club}: Mathematics outreach and lessons at local elementary school.\\
\end{tabular}

\section{\sc Talks}

\emph{Singular spaces in Poisson and symplectic geometry}, University of Toledo Department of Mathematics and Statistics Colloquium, Toledo, OH (October 2025).

\emph{Integration and differentiation for diffeological groupoids}, 5-day workshop on Poisson Geometry, Lie Groupoids and Differentiable Stacks. BIRS, Banff (June 2022).\\
Recording: \url{https://www.birs.ca/events/2022/5-day-workshops/22w5035/videos/watch/202206101030-Villatoro.html}

\emph{Diffeological Solutions to the Integration Problem} Gone Fishing (2020-2022). Savannah, GA (March 2022).

\emph{Diffeological vector bundles vs sheaves of modules}, Western University. London, Ontario CA. (October 2020).

\emph{On the homotopy theory and algebraic geometry of sheaves of Lie-Rinehart algebras}, Junior Global Poisson Workshop 2020. Online, (September 2020). \\
Recording: \url{https://www.youtube.com/watch?v=VFVr4xPkDRw&feature=emb_title}

\emph{A path approach to holonomy for singular foliations}, University of Illinois and Urbana-Champaign. Urbana, Illinois USA. (December 2019).

\emph{Holonomy for Singular Foliations}, University of Antwerp. Antwerp, Belgium. (November 2019).

\emph{Symplectic Morita Equivalence and Stacks}, Max Planck Institute for Mathematics. Bonn, Germany (February 2019).

\emph{On the singularities of the Weinstein groupoid}, Higher Geometric Structures along the Lower Rhine XII. Radboud University, Nijmegen, NL (January 2019).

\emph{Transversely Integrable Lie Algebroids and Asymmetric Orbispaces}, Utrecht University. Utrecht, Netherlands (October 2018).

\emph{Geometric stacks for groupoids with geometry}, Poisson 2018. University of Toronto, Toronto Canada (August 2018).

\emph{Picard groups of b-symplectic manifolds}, 3rd Pacific Rim Mathematical Association (PRIMA) congress. Instituto Tecnol\'ogico de Oaxaca, Oaxaca Mexico (August 2017).

\emph{Morita type equivalences for non-integrable algebroids}, Conference on Poisson Geometry and Stacks. The Fields Institute, University of Toronto (July 2017).

\emph{Poisson manifolds and their associated stacks}, 5-day workshop on Geometric Structures on Lie Groupoids. BIRS, Banff (April 2017).\\
Recording: \url{https://www.youtube.com/watch?v=rEzlq8118CQ}

\emph{Poisson manifolds and their stacks}, Symplectic and Poisson Geometry Seminar. University of Illinois at Urbana-Champaign (March 2017).

\emph{Picard Groups of b-Symplectic Manifolds} Symplectic and Poisson Geometry Seminar. University of Illinois at Urbana-Champaign (October 2016).

\emph{Motivating stacks as models for singular spaces}, Graduate Geometry and Topology seminar. University of Illinois at Urbana-Champaign (September 2016).

\emph{Poisson manifolds as families of Hamiltonian systems}, Graduate Geometry and Topology seminar. University of Illinois at Urbana-Champaign (September 2015).



\end{resume}
\end{document}
